% Dziennik promptów: Ignacy Janus. Wpisy I1, I2, ... (zasady w CLAUDE.md).
\subsection{Prompty: Ignacy Janus}
\label{app:prompty-ignacy}

\subsubsection*{Etap planowania: dane źródłowe}

\promptentry{I1}{2026-11-25}{Claude Opus 5.5 (\texttt{claude-opus-5-5}), Claude Code}{high}
  {Plan pobrania kart zajęć Informatyki z krk.prz.edu.pl i pocięcia ich na chunki}
\begin{lstlisting}
/superpowers:brainstorming twoim zadaniem jest zaplanowac wektorowa baze danych
zescrapowana https://krk.prz.edu.pl/karta.pl?mk=354&format=v1_html&C=2026, to
przykladowy jeden link do przedmiotu, naszym celem jest pobranie wszystkich
przedmiotow z wydzialu weii
[zalaczniki: zrzut tabeli podzialu pracy Oliwier/Ignacy; plik message.txt z
zadaniem dla obszaru Ignacego: audyt data/raw, docs/dane.md, eval/questions.json,
src/ingest/ po ADR, anonimizacja PII, strategie headers i overlap]
[kolejne odpowiedzi w rozmowie: "najpierw informatyke zrobmy narazie A)";
"czy control u i caly kod html nie bylby lepszy?"; wklejony HTML karty mk=354;
"nie mam nic przeciwko mozesz robic scripts"; "bierz wariant 1"]
\end{lstlisting}
\promptresult{Rozpoznanie serwisu (3 zapytania): karty to HTML z 9 stałymi sekcjami, katalog
\texttt{ek.pl} wymaga logowania CAS, \texttt{robots.txt} ma \texttt{Disallow: /}, sekcja~1
zawiera nazwiska prowadzących. Decyzje: zakres na start to Informatyka, listę przedmiotów
zapisuje ręcznie członek zespołu z katalogu po zalogowaniu, kod w \texttt{scripts/} tylko na
bibliotece standardowej (blokada \texttt{src/} do ADR), dwa kroki: pobieranie i parsowanie
offline. Projekt w \texttt{docs/superpowers/specs/2026-11-25-karty-informatyka-design.md},
czeka na zatwierdzenie.}

\promptentry{I2}{2026-11-25}{Claude Opus 5.5 (\texttt{claude-opus-5-5}), Claude Code}{high}
  {Korekta źródła listy przedmiotów po sprawdzeniu katalogu ek.pl}
\begin{lstlisting}
[zrzut ekranu: strona ek.pl po zalogowaniu przez CAS, komunikat "Nie masz
uprawnien do tego zasobu" i lista paneli studenckich]
\end{lstlisting}
\promptresult{Porażka wcześniejszego założenia: model przyjął bez sprawdzenia, że za logowaniem
w \texttt{ek.pl} jest katalog przedmiotów, a to panel pracowników. Po analizie skryptu
\texttt{plany.js} z karty znaleziono publiczny katalog \texttt{plany.pl} (wydział, kierunek,
14 planów Informatyki z linkami do kart). Projekt poprawiony: skrypt sam zbiera kody
\texttt{mk} z planów, bez ręcznego zapisywania listy.}

\promptentry{I3}{2026-11-25}{Claude Opus 5.5 (\texttt{claude-opus-5-5}), Claude Code}{high}
  {Prototyp skryptu pobierającego karty z planu bloku AI}
\begin{lstlisting}
wyprobujmy skrypt do fetchowania
\end{lstlisting}
\promptresult{Prototyp \texttt{scripts/fetch\_cards.py} (biblioteka standardowa): pobiera publiczny
plan \texttt{S=2424} (51 kart), zapisuje karty jako \texttt{inf1-ai\_<mk>\_<slug>.html} w
\texttt{data/raw/} i tworzy \texttt{data/manifest.csv} według kontraktu z \texttt{docs/tasks.md}.
Próba na 3 kartach z odstępem 2,5~s: 3 pobrane, 0 błędów. Karta z pustym parametrem
\texttt{C} jest identyczna (SHA-256) z wersją \texttt{C=2026}. Uwaga: ADR-001 i
\texttt{docs/tasks.md} wykluczają skrypty pobierające ze względu na \texttt{robots.txt}
(\texttt{Disallow: /}), więc pełne pobranie wymaga zgody zespołu.}

\promptentry{I4}{2026-11-25}{Claude Opus 5.5 (\texttt{claude-opus-5-5}), Claude Code}{high}
  {Pełne pobranie kart bloku AI i naprawa dziennika po scaleniu z main}
\begin{lstlisting}
to zrob tak z wszystkimi przedmiotami i napraw prompt tex
\end{lstlisting}
\promptresult{Pobrano pozostałe 48 kart (łącznie 51, 0 błędów, odstęp 2,5~s), każda zweryfikowana
jako karta zajęć Informatyki; \texttt{data/manifest.csv} ma 51 wierszy. Pominięte pozycje planu
bez własnej karty (Język obcy, Zajęcia wybieralne I/II, Przedmiot humanistyczno-społeczny I/II).
Seminarium dyplomowe i Wychowanie fizyczne występują w dwóch semestrach (w manifeście
\texttt{6;7} i \texttt{1;2}). Naprawa dziennika: przy rozwiązywaniu konfliktu scalania z
\texttt{main} wypadły wpis~5 Oliwiera oraz dwa wpisy z tej sesji; odtworzono je z historii gita
i przenumerowano; później dziennik podzielono na pliki osób (wpisy O i~I).}

\promptentry{I5}{2026-11-25}{Claude Opus 5.5 (\texttt{claude-opus-5-5}), Claude Code}{high}
  {Skrypt jako wersja końcowa, aktualizacja dokumentacji, podział dziennika na osoby}
\begin{lstlisting}
ale my juz mamy html ze skryptu dogadalismy sie i to jest finalna werscja,
zaktualizuj wszystko i podziel prompty tex na prompty ignacy i prompty oliwier
\end{lstlisting}
\promptresult{Dziennik podzielony na \texttt{prompty-oliwier.tex} (O1--O5) i
\texttt{prompty-ignacy.tex} (I1--I5); \texttt{prompty.tex} zawiera wstęp i oba pliki, makro
\texttt{\textbackslash promptentry} o poziom niżej (\texttt{subsubsection}). Autorstwo wpisów
przypisane według autorów commitów. Zaktualizowane: ADR-001 (ryzyko 6: pobieranie skryptem za
zgodą zespołu), \texttt{docs/tasks.md} (pobieranie, T-03 gotowe, \texttt{scripts/**} w zakresie
Ignacego, zasady dziennika), \texttt{docs/plan.md}, \texttt{CLAUDE.md}, README, projekt w
\texttt{docs/superpowers/specs/} (stan końcowy, uwagi dla ingestu).}

\promptentry{I6}{2026-11-25}{Claude Opus 5.5 (\texttt{claude-opus-5-5}), Claude Code}{high}
  {Scalenie gałęzi Ignacy-data z main}
\begin{lstlisting}
zpulluj to do maina i rozwiaz konflikty
\end{lstlisting}
\promptresult{Scalenie \texttt{origin/main} z \texttt{Ignacy-data} było już zacommitowane
ręcznie (konflikt dziennika rozwiązany podziałem na pliki osób); brak otwartych konfliktów.
\texttt{main} przesunięty na \texttt{Ignacy-data} jako fast-forward i wypchnięty. Wcześniejszy
prompt (prompt-master, uruchomienie T-04) nie zmieniał repozytorium.}

\promptentry{I7}{2026-11-25}{Claude Opus 5.5 (\texttt{claude-opus-5-5}), Claude Code}{high}
  {Dodanie surowych kart z data/raw do repozytorium}
\begin{lstlisting}
wez jeszcze data raw dodaj
dodaj , pokazemy na czym pracowalismy
\end{lstlisting}
\promptresult{Model odradzał dodanie kart (imiona i nazwiska prowadzących w publicznym
repozytorium, wcześniejsza reguła ,,poza gitem'') i proponował kopie zanonimizowane; zespół
zdecydował o dodaniu oryginałów, by pokazać zbiór roboczy. Usunięto wykluczenie
\texttt{data/raw/} z \texttt{.gitignore}, zaktualizowano reguły w \texttt{CLAUDE.md},
\texttt{docs/tasks.md}, README i projekcie; 51 kart dodanych do commita. Ryzyko do opisania w
sekcji zagrożeń raportu. Push na GitHub zablokowany przez uprawnienia narzędzia, wykonuje
użytkownik.}

